\documentclass[12pt,a4paper]{article}
\usepackage[utf8]{inputenc}
\usepackage{amsmath, amssymb, amsthm}
\usepackage{geometry}
\usepackage{xcolor}
\usepackage{physics}
\usepackage{siunitx}
\usepackage{upgreek}
\usepackage{bm}
\usepackage[version=4]{mhchem}
\usepackage{enumitem}
\usepackage{booktabs}
\usepackage{array}
\usepackage{hyperref}
\hypersetup{colorlinks=true, linkcolor=blue!60!black, citecolor=blue!60!black}

\geometry{top=2.5cm, bottom=2.5cm, left=2.5cm, right=2.5cm}

\newcommand{\R}{\mathbb{R}}
\newcommand{\oN}{\mathbf{1}_N}
\newcommand{\oFour}{\mathbf{1}_4}
% Dimension symbols (upright sans-serif), distinct from variables
\newcommand{\dM}{\mathsf{M}}
\newcommand{\dL}{\mathsf{L}}
\newcommand{\dT}{\mathsf{T}}
\newcommand{\Cw}{C_{\mathrm{w}}}

\title{\textbf{Mechanistic Compartmental Model for PFAS Bioaccumulation in Rice}\\[6pt]
\large An Ionizable-Organic-Compound (IOC) Extension of the Dynamic Plant Uptake Framework}
\author{Chan Young Joe}
\date{\today}

\begin{document}
\maketitle
\tableofcontents
\newpage

%% ============================================================
\section{Introduction and Scope}
%% ============================================================

This report develops a mechanistic four-compartment model for the uptake and tissue distribution of per- and polyfluoroalkyl substances (PFAS) in rice (\emph{Oryza sativa}).  It extends the neutral-compound Dynamic Plant Uptake (DPU) base framework to \textbf{ionizable organic compounds (IOC)}, with PFAS treated as the limiting case of a \emph{permanently dissociated anion} (very low $\mathrm{p}K_\mathrm{a}$, dissociated fraction $f_d\approx1$ at environmental pH).

Three features distinguish PFAS-in-rice from the neutral DPU base and motivate every structural change below:
\begin{enumerate}[nosep]
    \item \textbf{Charge.} The octanol--water / Briggs partition core of the neutral model does not apply; for PFAS the translocation factor correlates \emph{negatively} with $\log K_\mathrm{OW}$, and $K_\mathrm{OW}$ is itself ill-defined for surfactants.  Membrane crossing is governed by electrodiffusion together with channel- and carrier-mediated transport.
    \item \textbf{Phloem-fed grain.} The rice grain is a terminal sink filled through the phloem, which the xylem-only DPU base cannot represent.  Moreover, the classical weak-acid \emph{ion-trap} phloem mechanism does not operate for a fully dissociated anion ($f_n\approx0$), so phloem loading is recast as carrier/channel loading.
    \item \textbf{Binding-driven accumulation.} Tissue retention is set by protein, phospholipid, and cell-wall binding rather than by neutral-lipid partitioning; this is encoded in compartment binding factors $B_k$.
\end{enumerate}

\paragraph{Convention.}  Dimension symbols are upright sans-serif ($\dM$ mass, $\dL$ length, $\dT$ time, $\mathrm{mol}$ amount).  A dot marks a rate (a flow or a flux contributing to $\dv*{C}{t}$).  The state variable in each compartment $k$ is the total tissue concentration $C_k\;[\mathrm{mol\,kg^{-1}}]$ on a fresh-mass basis.


%% ============================================================
\section{Model Structure and Assumptions}
%% ============================================================

The plant is represented by four compartments: roots ($k=1$), stem ($k=2$), leaves ($k=3$), and fruits/grain ($k=4$).  The core assumptions are:

\begin{description}[leftmargin=2.2em,style=nextline,nosep]
    \item[(A1)] Internal compartments exchange solute by \emph{advection} (xylem and phloem) plus internal partitioning; an explicit transmembrane uptake flux $j_R$ is retained only at the soil--root boundary.  The cell-level electrodiffusion equilibrium for internal tissues is absorbed into the binding factors $B_k$.
    \item[(A2)] Each compartment is well mixed; the concentration carried by the sap leaving a compartment is the free aqueous concentration $\Cw{}_{,k}=C_k/B_k$, except at the root--xylem boundary, where only a fraction $f_\mathrm{xy}\le1$ (the transpiration-stream concentration factor, TSCF) is loaded into the ascending xylem.  $f_\mathrm{xy}\to1$ recovers the unrestricted DPU base; $f_\mathrm{xy}\ll1$ for anions, whose endodermal (Casparian) loading and strong root binding sequester the solute in the root and limit root-to-shoot translocation.
    \item[(A3)] PFAS are non-volatile ($K_\mathrm{AW}\approx0$); volatilisation, gaseous uptake, and particle deposition are neglected.
    \item[(A4)] PFAS are metabolically recalcitrant; first-order metabolisation rate constants $\gamma_k\approx0$ (kept symbolically, set to zero in the PFAS limit).
    \item[(A5)] PFAS are fully dissociated ($f_d\approx1$, $f_n\approx0$); the neutral-permeation pathway and the weak-acid pH ion-trap are inactive.
    \item[(A6)] Tissue mass $M_k(t)$, transpiration $\dot Q_\mathrm{TP}(t)$, and the soil pore-water free concentration $C_\mathrm{w}^{o}(t)$ are external inputs.
\end{description}


%% ============================================================
\section{Root Membrane Uptake: Hybrid Flux}\label{sec:root}
%% ============================================================

\subsection{General hybrid flux}

Across the root plasma membrane, the molar flux per unit membrane area $J_R\;[\mathrm{mol\,m^{-2}\,s^{-1}}]$ between soil pore water (superscript $o$) and the root symplast aqueous phase is the parallel sum of three pathways:
\begin{equation}\label{eq:JR}
    J_R=\underbrace{P_n\!\left(f_n^{o}\Cw^{o}-f_n^{i}\Cw{}_{,1}\right)}_{\text{neutral passive}}
    +\underbrace{P_d^{\mathrm{eff}}\,\frac{N}{e^{N}-1}\!\left(f_d^{o}\Cw^{o}-f_d^{i}e^{N}\Cw{}_{,1}\right)}_{\text{ionic electrodiffusion}}
    +\underbrace{\frac{V_{\max}^{\mathrm{in}}\Cw^{o}}{K_m^{\mathrm{in}}+\Cw^{o}}-\frac{V_{\max}^{\mathrm{out}}\Cw{}_{,1}}{K_m^{\mathrm{out}}+\Cw{}_{,1}}}_{\text{carrier-mediated (saturable)}},
\end{equation}
where $N=zEF/(RT)$ is the dimensionless electrochemical driving force ($z=-1$ for PFAS, $E$ the membrane potential, $E\approx-\SI{120}{mV}$), and $\Cw{}_{,1}=C_1/B_1$ is the root free aqueous concentration.  The passive channels (anion channels, aquaporins) act in parallel with the lipid bilayer and are absorbed into the effective ionic permeability,
\begin{equation}\label{eq:Pdeff}
    P_d^{\mathrm{eff}}=P_d+P_\mathrm{ac}+P_\mathrm{aq},\qquad P_d\approx P_n\cdot10^{-3.5},
\end{equation}
with $P_\mathrm{ac}$ (anion channel) and $P_\mathrm{aq}$ (aquaporin) as conductive contributions.  Only the carrier term ($V_{\max}$) is energy-dependent, i.e.\ the metabolic-inhibitor-sensitive fraction.  The mass-specific root uptake is $j_R=a_R J_R$, with $a_R\;[\mathrm{m^2\,kg^{-1}}]$ the specific root membrane area.

\subsection{PFAS limit}

For a fully dissociated anion ($f_n\to0,\;f_d\to1$) the neutral term vanishes:
\begin{equation}\label{eq:JR_pfas}
    J_R\approx P_d^{\mathrm{eff}}\,\frac{N}{e^{N}-1}\!\left(\Cw^{o}-e^{N}\Cw{}_{,1}\right)
    +\frac{V_{\max}^{\mathrm{in}}\Cw^{o}}{K_m^{\mathrm{in}}+\Cw^{o}}-\frac{V_{\max}^{\mathrm{out}}\Cw{}_{,1}}{K_m^{\mathrm{out}}+\Cw{}_{,1}}.
\end{equation}
The electrodiffusion term alone gives, at $J=0$, $\Cw{}_{,1}/\Cw^{o}=e^{-N}=e^{EF/RT}<1$: passive electrodiffusion \emph{excludes} the anion (inside-negative membrane).  Observed accumulation must therefore be driven by the carrier influx $V_{\max}^{\mathrm{in}}$ and/or by binding (Section~\ref{sec:binding}).  This structure naturally separates congeners: PFOA behaves carrier-dominated (active), PFOS electrodiffusion/aquaporin-dominated (passive).

\subsection{Free versus total concentration: binding}\label{sec:binding}

The compartment state variable is the total tissue concentration, related to the free aqueous concentration through a (linearised) binding capacity factor:
\begin{equation}\label{eq:binding}
    C_k=B_k\,\Cw{}_{,k},\qquad
    B_k=\theta_k+f_\mathrm{prot}K_\mathrm{prot}+f_\mathrm{PL}K_\mathrm{PL}+f_\mathrm{cw}K_\mathrm{cw},
\end{equation}
with $\theta_k$ the aqueous volume fraction $[\mathrm{L\,kg^{-1}}]$ and $f_\mathrm{prot},f_\mathrm{PL},f_\mathrm{cw}$ the protein, phospholipid, and cell-wall mass fractions $[\mathrm{kg\,kg^{-1}}]$ multiplying the partition coefficients $K_i\,[\mathrm{L\,kg^{-1}}]$.  There is \emph{no} density prefactor: $B_k$ is a per-fresh-mass capacity, so each term already carries units $\mathrm{L\,kg^{-1}}$ and $C_k=B_k\Cw{}_{,k}$ is dimensionally consistent.  Protein binding saturates at high concentration; $K_\mathrm{prot}$ may be replaced by a Langmuir form where needed.  $B_k$ carries the binding-driven accumulation; the large protein content of grain makes $B_4$ large, trapping PFAS in the edible tissue.


%% ============================================================
\section{Governing Equations for the Four-Compartment System}\label{sec:system}
%% ============================================================

\subsection{Long-distance transport}

The xylem flow follows the DPU distribution
\begin{equation}\label{eq:Qxyl}
    \dot{\mathbf Q}_\mathrm{XYL}=\Big[\dot Q_\mathrm{TP},\;\dot Q_\mathrm{TP},\;\tfrac{A_3}{A_3+A_4}\dot Q_\mathrm{TP},\;\tfrac{A_4}{A_3+A_4}\dot Q_\mathrm{TP}\Big]^{\mathsf T},
\end{equation}
and each compartment exports sap at its free concentration $C_k/B_k$, with the exception of the root, whose xylem loading is throttled by the transpiration-stream concentration factor $f_\mathrm{xy}\le1$ (A2): the concentration entering the ascending xylem is
\begin{equation}\label{eq:Cxyl}
    \Cw^{\mathrm{xyl}}=f_\mathrm{xy}\,\frac{C_1}{B_1}.
\end{equation}
The phloem flow is driven by grain sink demand plus a recirculating fraction $\varphi$ of transpiration:
\begin{equation}\label{eq:Qphl}
    \dot Q_\mathrm{Phl}=\dv{M_4}{t}\,T_{C,\mathrm{Ph}}+\varphi\,\dot Q_\mathrm{TP},
\end{equation}
with phloem sap concentration set by carrier/channel loading at the leaf source (\emph{not} a pH ion-trap, per A5):
\begin{equation}\label{eq:Cphl}
    C_\mathrm{Phl}=L_\mathrm{Ph}\,\frac{C_3}{B_3}.
\end{equation}

\subsection{Compartment mass balances}

Writing $\mu_k=\tfrac{1}{M_k}\dv*{M_k}{t}$ for the growth-dilution rate,
\begin{align}
    \dv{C_1}{t}&=j_R-\frac{\dot Q_\mathrm{TP}}{M_1}\Cw^{\mathrm{xyl}}+\varphi\frac{\dot Q_\mathrm{Phl}}{M_1}C_\mathrm{Phl}-\gamma_1C_1-\mu_1C_1, \label{eq:root}\\[4pt]
    \dv{C_2}{t}&=\frac{\dot Q_\mathrm{TP}}{M_2}\!\left(\Cw^{\mathrm{xyl}}-\frac{C_2}{B_2}\right)-\gamma_2C_2-\mu_2C_2, \label{eq:stem}\\[4pt]
    \dv{C_3}{t}&=\frac{A_3}{A_3+A_4}\frac{\dot Q_\mathrm{TP}}{M_3}\frac{C_2}{B_2}-(1+\varphi)\frac{\dot Q_\mathrm{Phl}}{M_3}C_\mathrm{Phl}-\gamma_3C_3-\mu_3C_3, \label{eq:leaf}\\[4pt]
    \dv{C_4}{t}&=\frac{A_4}{A_3+A_4}\frac{\dot Q_\mathrm{TP}}{M_4}\frac{C_2}{B_2}+\frac{\dot Q_\mathrm{Phl}}{M_4}C_\mathrm{Phl}-\gamma_4C_4-\mu_4C_4. \label{eq:fruit}
\end{align}
With $j_R=a_RJ_R$ from Eq.~\eqref{eq:JR_pfas} and the closures \eqref{eq:binding} and \eqref{eq:Cxyl}--\eqref{eq:Cphl}, Eqs.~\eqref{eq:root}--\eqref{eq:fruit} form a closed nonlinear ODE system in $(C_1,\dots,C_4)$.  As the phloem source, the leaf loses the \emph{whole} phloem export $(1+\varphi)\dot Q_\mathrm{Phl}C_\mathrm{Phl}$, which feeds the grain sink ($\dot Q_\mathrm{Phl}C_\mathrm{Phl}$) and the root recirculation ($\varphi\dot Q_\mathrm{Phl}C_\mathrm{Phl}$); together with the matched root$\to$stem loading $\Cw^{\mathrm{xyl}}$, the internal xylem and phloem transfers then conserve solute mass exactly, leaving $j_R$ as the only net source (the only sink being metabolism, $\gamma_k\to0$).


%% ============================================================
\section{Nondimensionalization and Dimensionless Groups}\label{sec:nondim}
%% ============================================================

Scaling free concentrations by the soil value, $u_k=\Cw{}_{,k}/\Cw^{o}=C_k/(B_k\Cw^{o})$, the system is governed by:
\begin{center}
\renewcommand{\arraystretch}{1.3}
\begin{tabular}{@{}ll@{}}
\toprule
\textbf{Group} & \textbf{Meaning} \\
\midrule
$B_k$ & binding capacity (tissue retention vs re-export) \\
$N=zEF/RT$ & electrochemical driving force \\
$f_\mathrm{xy}$ & root$\to$xylem loading factor (TSCF; translocation throttle) \\
$\mathrm{Da}_k=\gamma_k B_k M_k/\dot Q_\mathrm{TP}$ & metabolism vs advective throughput \\
$\Pi=\dot Q_\mathrm{Phl}L_\mathrm{Ph}/\dot Q_\mathrm{TP}$ & phloem vs xylem (leaf$\to$grain partitioning) \\
$g_\mathrm{in}/g_\mathrm{out}$ & root influx vs efflux conductance (Eq.~\ref{eq:gio}) \\
\bottomrule
\end{tabular}
\end{center}
For PFAS, $\gamma_k\approx0\Rightarrow\mathrm{Da}_k\approx0$, a substantial simplification.


%% ============================================================
\section{Quasi-Steady-State Analysis}\label{sec:qss}
%% ============================================================

Assuming the chemical distribution equilibrates fast relative to growth ($\dv*{C_k}{t}=0$, $\gamma_k\approx0$):

\paragraph{Xylem conservation.}  From Eq.~\eqref{eq:stem},
\begin{equation}
    \frac{u_2}{u_1}=\frac{f_\mathrm{xy}}{1+\mathrm{Da}_2}\;\xrightarrow{\ \gamma_2\to0\ }\;f_\mathrm{xy},
\end{equation}
so for non-metabolised PFAS the free concentration entering the shoot is the root value reduced by the loading factor ($\Cw{}_{,2}\approx f_\mathrm{xy}\Cw{}_{,1}$) and is then conserved along the xylem.

\paragraph{Root bioaccumulation factor.}  In the low-concentration (linear) limit $\Cw\ll K_m$, define
\begin{equation}\label{eq:gio}
    g_\mathrm{in}=a_R\!\left[\frac{P_d^{\mathrm{eff}}N}{e^{N}-1}+\frac{V_{\max}^{\mathrm{in}}}{K_m^{\mathrm{in}}}\right],\qquad
    g_\mathrm{out}=a_R\!\left[\frac{P_d^{\mathrm{eff}}N}{e^{N}-1}e^{N}+\frac{V_{\max}^{\mathrm{out}}}{K_m^{\mathrm{out}}}\right].
\end{equation}
Neglecting the (small) phloem return, the root balance \eqref{eq:root} at steady state gives
\begin{equation}\label{eq:BAF}
    \mathrm{BAF}_1=\frac{C_1}{\Cw^{o}}=B_1\,\frac{g_\mathrm{in}}{\,g_\mathrm{out}+f_\mathrm{xy}\dot Q_\mathrm{TP}/M_1\,}.
\end{equation}
The root BAF is thus the binding capacity $B_1$ times an influx/(efflux $+$ loaded xylem export) ratio.  A small loading factor $f_\mathrm{xy}$ both raises the root BAF (less is exported upward) and starves the shoot, which is the mechanism that restores the empirical root~$>$~straw~$>$~grain ordering.  The chain-length dependence observed empirically (concave/U-shaped for PFCA, monotone for PFSA) enters through the chain-length variation of $B_1$, of $g_\mathrm{in}/g_\mathrm{out}$, and of $f_\mathrm{xy}$.

\paragraph{Grain as a terminal accumulator.}  The fruit \eqref{eq:fruit} has no outflow and $\gamma_4\approx0$; the only sink is growth dilution $\mu_4=K_4^\mathrm{gr}(1-M_4/M_4^\mathrm{max})$, which vanishes as the grain matures.  Hence \emph{no bounded steady state exists}: the grain is an accumulator, and its concentration at harvest is the time integral of all inflows divided by the final grain mass,
\begin{equation}\label{eq:grain}
    C_4^{\mathrm{harv}}=\frac{1}{M_4^{\mathrm f}}\int_0^{t_\mathrm{harv}}\!\left(\dot Q_\mathrm{Phl}\,C_\mathrm{Phl}+\frac{A_4}{A_3+A_4}\dot Q_\mathrm{TP}\,\frac{C_2}{B_2}\right)dt.
\end{equation}
This is precisely why a \emph{dynamic} model coupled to grain-filling phenology is required; a steady-state formulation cannot, in principle, predict the edible-tissue concentration.

\paragraph{Identifiability.}  Equation~\eqref{eq:BAF} depends on $P_d^{\mathrm{eff}}$ and $V_{\max}$ only through the \emph{lumped} conductance $g_\mathrm{in}$, not separately.  Bioaccumulation data alone therefore cannot resolve the channel (passive) from the carrier (active) contribution; selective-inhibitor experiments are required to separate them.


%% ============================================================
\section{Parameter Inventory and Calibration Design}\label{sec:params}
%% ============================================================

Parameters are organised by the data required to determine them.

\paragraph{Tier 0 --- inputs / known (fixed).}
\begin{center}
\renewcommand{\arraystretch}{1.25}
\begin{tabular}{@{}p{4.4cm}p{6.2cm}p{3.6cm}@{}}
\toprule
\textbf{Symbol} & \textbf{Meaning} & \textbf{Source} \\
\midrule
$M_k(t),\dot Q_\mathrm{TP}(t),\Cw^{o}(t)$ & mass, transpiration, soil free conc.\ (drivers) & measurement / sub-model \\
$M_k^0,M_k^\mathrm{max},K_k^\mathrm{gr}$ & logistic growth & biomass curves \\
$S_k,\rho_k,\theta_k$ & specific area, density, water fraction & tissue measurement \\
$N=zEF/RT$ & electrochemical number & $E\approx-\SI{120}{mV}$, $z=-1$ \\
$f_d\approx1,\,f_n\approx0$ & dissociation fractions & $\mathrm{p}K_\mathrm{a}$ (strong acid) \\
$\gamma_k\approx0$ & metabolism & recalcitrant $\to0$ \\
$T_{C,\mathrm{Ph}}$ & phloem flux coefficient & literature ($\sim\SI{10}{L/kg}$) \\
\bottomrule
\end{tabular}
\end{center}

\paragraph{Tier 1 --- identifiable from BAF / tissue-concentration data (lumped).}
\begin{center}
\renewcommand{\arraystretch}{1.25}
\begin{tabular}{@{}p{4.4cm}p{6.2cm}p{3.6cm}@{}}
\toprule
\textbf{Symbol} & \textbf{Meaning} & \textbf{Data} \\
\midrule
$B_k$ & compartment binding capacity & per-tissue BAF \\
$g_\mathrm{in}/g_\mathrm{out}$ & root influx/efflux conductance ratio & root BAF \\
$f_\mathrm{xy}$ & root$\to$xylem loading factor (TSCF) & root vs shoot BAF \\
$\Pi=\dot Q_\mathrm{Phl}L_\mathrm{Ph}/\dot Q_\mathrm{TP}$ & phloem/xylem ratio & shoot \& grain BAF \\
$\varphi$ & phloem recirculation fraction & root--shoot pattern \\
\bottomrule
\end{tabular}
\end{center}

\paragraph{Tier 2 --- require inhibitor / kinetic experiments to separate.}
\begin{center}
\renewcommand{\arraystretch}{1.25}
\begin{tabular}{@{}p{4.4cm}p{6.2cm}p{3.6cm}@{}}
\toprule
\textbf{Split} & \textbf{Meaning} & \textbf{Experiment} \\
\midrule
$P_d\mid P_\mathrm{ac}\mid P_\mathrm{aq}$ & membrane / anion-channel / aquaporin & channel \& aquaporin blockers \\
$V_{\max},K_m\mid P_d^{\mathrm{eff}}$ & carrier vs passive & metabolic inhibitor + conc.\ series \\
$V_{\max}^{\mathrm{in}}\mid V_{\max}^{\mathrm{out}}$ & influx/efflux asymmetry & uptake/depuration kinetics \\
\bottomrule
\end{tabular}
\end{center}

\paragraph{Tier 3 --- independent measurement / QSPR (chain-length resolved).}
\begin{center}
\renewcommand{\arraystretch}{1.25}
\begin{tabular}{@{}p{4.4cm}p{6.2cm}p{3.6cm}@{}}
\toprule
\textbf{Symbol} & \textbf{Meaning} & \textbf{Source} \\
\midrule
$K_\mathrm{prot}$ & PFAS--protein binding & assay (albumin/FABP) / QSPR \\
$K_\mathrm{PL}$ & phospholipid--water partition & measurement / QSPR \\
$K_\mathrm{cw}$ & cell-wall sorption & measurement / estimate \\
$L_\mathrm{Ph}$ & phloem loading partition & sparse data $\to$ estimate/calibrate \\
$a_R$ & specific root membrane area & anatomical measurement \\
\bottomrule
\end{tabular}
\end{center}

\paragraph{Calibration workflow.}  The tiers define the procedure: fix Tier~0; supply Tier~3 by QSPR/assay (decomposing $B_k$ and providing $L_\mathrm{Ph}$); calibrate Tier~1 against chain-length-resolved BAF data (chain-length trends can be shared across homologues through a hierarchical structure); and obtain Tier~2 from dedicated inhibitor and uptake/depuration experiments.  The dominant uncertainties are expected to be (i) $K_\mathrm{prot}$ (controls grain $B_4$ and edible-tissue accumulation), (ii) $L_\mathrm{Ph}$ (data-poor phloem loading), and (iii) the $P_d^{\mathrm{eff}}$ vs $V_{\max}$ separation (not identifiable from BAF alone).


%% ============================================================
\section{Soil Boundary Condition and Calibration (Method A)}
%% ============================================================

\paragraph{Paddy soil sorption.}  The driver $\Cw^{o}(t)$ is obtained from a soil PFAS inventory through a Freundlich partition between pore water and the solid phase,
\begin{equation}\label{eq:freundlich}
    S=K_F\,\Cw^{\,n},\qquad
    C_T=K_F\,\Cw^{\,n}+\theta_g\,\Cw,
\end{equation}
with $S$ the sorbed solid-phase load, $C_T$ the total PFAS per unit dry soil mass, $\theta_g$ the gravimetric water content $[\mathrm{L\,kg^{-1}}]$, $K_F$ the Freundlich capacity, and $n$ the (typically sub-linear) exponent.  Since $C_T(\Cw)$ is strictly increasing for $\Cw\ge0$, the bioavailable pore-water concentration $\Cw^{o}=\Cw$ is recovered by inverting Eq.~\eqref{eq:freundlich} (the linear $K_d$ limit is $n=1$).  Flooding versus draining the paddy switches the redox state and hence $(K_F,n)$, so $\Cw^{o}(t)$ follows the water-management schedule.  This sub-model is interchangeable with a HYDRUS-1D/Phydrus solute output under the loose, one-way Method~A coupling.

\paragraph{Calibration.}  The Tier-1 parameters $\bm\beta=(B_k,\,g_\mathrm{in}/g_\mathrm{out},\,f_\mathrm{xy},\,\Pi,\,\varphi)$ are estimated by weighted least squares in log space against observed tissue bioaccumulation factors,
\begin{equation}\label{eq:objective}
    \hat{\bm\beta}=\arg\min_{\bm\beta}\;\sum_{k}\frac{1}{\sigma_k^{2}}\Big(\log\mathrm{BAF}_k^{\mathrm{pred}}(\bm\beta)-\log\mathrm{BAF}_k^{\mathrm{obs}}\Big)^{2},
\end{equation}
under box constraints, with the mass-weighted shoot (straw) treated as one observable.  Consistent with the identifiability analysis of Section~\ref{sec:qss}, BAF data recover only the lumped influx conductance $g_\mathrm{in}$: the channel ($P_d^{\mathrm{eff}}$) and carrier ($V_{\max}$) contributions trade off at fixed $g_\mathrm{in}$ and require inhibitor experiments to separate.


%% ============================================================
\section{Outlook}
%% ============================================================

The model skeleton, the dimensionless structure, the paddy soil boundary condition (Eq.~\ref{eq:freundlich}), and the calibration design (Eq.~\ref{eq:objective}) are implemented in the accompanying Python module under the Method~A coupling.  Items still deferred to subsequent work include: the QSPR construction for $K_\mathrm{prot}$ and $K_\mathrm{PL}$ (Tier~3); the saturable (nonlinear) regime of the carrier and protein-binding terms; the parameter fit against measured rice field and lysimeter datasets; and the tight (Method~B) coupling inside HYDRUS-1D.

\end{document}
